\section*{Balanced upper-mass composition and exact parity inflation}

For a Boolean function $f$ write
\[
 p_f(S)=\widehat f(S)^2,\qquad
 \pi_f(S)=\sum_{T\supseteq S}p_f(T),\qquad
 U(f)=\sum_{S:p_f(S)>0}p_f(S)\log_2\frac{\pi_f(S)}{p_f(S)}.
\]
For a physical coordinate $i$ let
\[
 \Delta_iU(f)=U(f)-\frac{U(f_{i,-})+U(f_{i,+})}{2},
 \qquad G(f)=\sum_i\Delta_iU(f).
\]

\subsection*{Theorem 1 (balanced disjoint substitution)}

Let $g:\{-1,1\}^m\to\{-1,1\}$ and let
$h_j:\{-1,1\}^{B_j}\to\{-1,1\}$ be Boolean functions on pairwise
disjoint physical blocks.  Suppose $\mathbb E h_j=0$ for every $j$, and set
\[
 F(x)=g(h_1(x_{B_1}),\ldots,h_m(x_{B_m})).
\]
Then
\[
 \boxed{\ U(F)=U(g)+\sum_{j=1}^m\operatorname{Inf}_j(g)U(h_j)\ },
 \qquad
 \boxed{\ I(F)=\sum_{j=1}^m\operatorname{Inf}_j(g)I(h_j)\ }.
\]

\paragraph{Proof.}
For a physical Fourier set $S$, let $S_j=S\cap B_j$ and
$A=\{j:S_j\ne\varnothing\}$.  Balancedness gives
\[
 \widehat F(S)=\widehat g(A)\prod_{j\in A}\widehat h_j(S_j),
 \qquad
 p_F(S)=p_g(A)\prod_{j\in A}p_{h_j}(S_j).
\]
In the upper-set sum, every newly active block $j\in C\setminus A$ is
summed over all nonempty Fourier sets.  Its total mass is one because
$h_j$ is balanced.  Consequently
\[
 \pi_F(S)=\pi_g(A)\prod_{j\in A}\pi_{h_j}(S_j).
\]
The logarithm of $\pi_F(S)/p_F(S)$ is therefore the outer logarithm plus
the inner logarithms.  Averaging the outer term gives $U(g)$, while the
$j$th inner term has weight
$\sum_{A\ni j}p_g(A)=\operatorname{Inf}_j(g)$.  This proves the first
identity.  The second follows by weighting the same factorized law by
$|S|=\sum_{j\in A}|S_j|$. \qed

This calculation also locates the general obstruction: after restricting
a coordinate of a generic balanced child, the two child functions need not
remain balanced, so this identity cannot simply be subtracted at the two
restricted nodes.

\subsection*{Theorem 2 (exact parity-inflation law)}

Let $B_1,\ldots,B_m$ be disjoint nonempty blocks of sizes
$d_1,\ldots,d_m$, put $P_j(x)=\prod_{r\in B_j}x_r$, and define
\[
 F(x)=g(P_1(x),\ldots,P_m(x)).
\]
For arbitrary Boolean $g$ and arbitrary positive integers $d_j$,
\[
 \boxed{\ U(F)=U(g)\ },\qquad
 \boxed{\ I(F)=\sum_{j=1}^m d_j\operatorname{Inf}_j(g)\ },
\]
and, more sharply, every $r\in B_j$ satisfies
\[
 \boxed{
 \Delta_rU(F)=
 \begin{cases}
 \Delta_jU(g),&d_j=1,\\
 0,&d_j\ge2.
 \end{cases}}
\]
Thus
\[
 \boxed{\ G(F)=\sum_{j:d_j=1}\Delta_jU(g)\ }.                 \tag{1}
\]

\paragraph{Proof.}
Each $P_j$ is balanced and has $U(P_j)=0$, so Theorem 1 gives $U(F)=U(g)$.
Equivalently, the Fourier coefficient $\widehat g(A)$ is merely moved to
the physical set $\bigcup_{j\in A}B_j$; the upper-set ratios on these
occupied atoms are unchanged.  Their physical degrees are
$\sum_{j\in A}d_j$, proving the influence formula.

Fix $r\in B_j$.  If $d_j\ge2$, restriction of $r$ replaces $P_j$ by a
possibly signed nonempty parity on $B_j\setminus\{r\}$.  The sign is an
input flip of $g$, which changes Fourier signs but not squared coefficients.
All parity blocks are still nonempty, so each restricted child has upper
mass exactly $U(g)$.  Hence $\Delta_rU(F)=0$.  If $d_j=1$, restriction of
$r$ sets the $j$th outer input.  The remaining nonempty parity substitution
preserves upper mass, so its two values are $U(g_{j,-})$ and $U(g_{j,+})$.
This gives $\Delta_rU(F)=\Delta_jU(g)$ and proves (1). \qed

\paragraph{Dimension-free covered class.}
If every $d_j\ge2$, then
\[
 G(F)=0\le 0\cdot I(F).
\]
This holds for every outer Boolean gate, every arity, every block size, and
every nested parity realization after flattening.  The constant is exactly
zero and has no dependence on dimension, arity, block size, or depth.  It is
a proper class result for the GU target; it does not settle universal GU or
FEI.

\subsection*{Proposition 3 (signed-decrement obstruction)}

Coordinate decrements $\Delta_jU(g)$ need not be nonnegative.  In truth-table
order
\[
 (-,-,-),(+,-,-),(-,+,-),(+,+,-),(-,-,+),(+,-,+),(-,+,+),(+,+,+),
\]
take
\[
 g=(+,+,-,+,+,-,-,-).
\]
Its Walsh coefficients, in subset-mask order, are
\[
 (0,0,\tfrac12,\tfrac12,\tfrac12,-\tfrac12,0,0).
\]
The two restrictions in coordinate $1$ are characters and have $U=0$.
The two restrictions in each of coordinates $2,3$ have four spectral atoms
of mass $1/4$ and have $U=1$.  Directly, $U(g)=1/2$, whence
\[
 (\Delta_1U,\Delta_2U,\Delta_3U)
   =(\tfrac12,-\tfrac12,-\tfrac12),
 \qquad G(g)=-\tfrac12,
 \qquad I(g)=\tfrac32.
\]
Inflating only the last two coordinates to parity blocks of any sizes at
least two deletes the two negative summands in (1).  In particular widths
$(1,2,2)$ give
\[
 U(F)=\tfrac12,\qquad G(F)=\tfrac12,\qquad I(F)=\tfrac52.
\]
Therefore neither coordinatewise nonnegativity of $\Delta_iU$ nor monotonicity
of $G$ under parity inflation is available.  Any GU composition proof that
drops or widens blocks must retain the signed coordinate ledger.  This is an
exact genuine-Boolean obstruction, not a counterexample to GU or FEI.

\subsection*{Verification scope and remaining gap}

The companion checker uses exact rational Walsh coefficients and 100-digit
decimal logarithms.  It verified Theorem 1 on all $16$ outer two-bit Boolean
gates and all $8^2$ ordered pairs of balanced one- or two-bit children
($1024$ checks).  It verified Theorem 2 for every Boolean outer gate of arity
at most three and every width vector in $\{1,2\}^m$ ($2120$ checks), including
$276$ all-nonterminal cases with $G=0$.

The unpaid universal step is the generic biased-child defect: for a balanced
child $h_j$ whose physical restriction is not another balanced character,
one must compare the two biased upper-mass laws and pay their signed coherent
remainder by a common influence or terminal-reserve budget.  The exact
balanced root chain rule alone does not perform this payment.
